单细胞转录组聚类需要同时解释细胞分组及其关联基因。本文复现CAdir基于对应分析与方向聚类的计算流程，在Tabula Muris Limb Muscle（LM）数据上实现原论文Fig3 A–D和关联S1 A–E，并在Nestorowa造血干／祖细胞（HSC/HSPC）数据上开展新数据对照实验。LM输入包含\allowbreak{}13204个基因和\allowbreak{}1882个细胞，使用既有对数归一化表达，选择\allowbreak{}6000个高变基因，计算\allowbreak{}30维对应分析，以初始簇数\allowbreak{}8、方向角阈值\allowbreak{}55°和随机种子\allowbreak{}2358执行CAdir。结果得到\allowbreak{}7个细胞簇及\allowbreak{}1200个已分配基因，调整兰德指数为\allowbreak{}0.669519，最大熵归一化互信息为\allowbreak{}0.741885。细胞数、已分配基因数、最终簇数、全部逐簇汇总及可见拆分合并日志与作者公开运行记录一致，调整兰德指数在作者显示精度下一致。九项分析包括关联图、方向比较、拆分合并历史、方向角矩阵、参考标签流量及坐标转换。新数据实验采用统一质量控制、归一化和\allowbreak{}6000个高变基因，保留\allowbreak{}1656个细胞，在\allowbreak{}1641个具有参考标签的细胞上评价CAdir和主成分分析＋共享近邻图＋Louvain两种流程。两种方法各运行五个预定种子，全部成功。CAdir的调整兰德指数中位数为\allowbreak{}0.291506，基线为\allowbreak{}0.269345；归一化互信息中位数分别为\allowbreak{}0.343414和\allowbreak{}0.356464。结果表明，在本次HSC配置下，两种评价指标呈现不同的相对排序；CAdir还提供了细胞与基因的联合分组、基因排名和方向诊断。本文结合数学表达、实际代码、运行命令、真实图表和独立数值复核，对上述结果进行解释。
\par
